\section{Conclusion}\label{sec:conclusion}

\subsection{Multi Criteria Analysis}
The multi-criteria analysis shows that Flutter, React Native, and Ionic all meet the essential communication requirements.
Each framework supports RESTful API calls and WebSocket connections.

The largest differences occur in the Community and Ecosystem criteria.
React Native achieves the highest Community score of \emph{0.93}, due to its large number of GitHub examples and its extensive surrounding community.
Flutter follows with a score of \emph{0.80}, supported by the highest number of GitHub stars and a large number of examples.
Ionic receives the lowest Community score of \emph{0.21}, as it has fewer stars, Discord users, and GitHub examples than the other frameworks.

Developer Experience scores are close.
React Native receives the highest score of \emph{0.74}, followed by Flutter with \emph{0.73} and Ionic with \emph{0.71}.
All frameworks provide debugging support and sufficient documentation.
The small differences are mainly caused by the team's prior experience, with React Native and Ionic benefiting from transferable JavaScript knowledge.

React Native and Ionic receive the maximum Ecosystem score of \emph{1.00}.
Both use npm, which contains more than two million packages, and provide suitable UI component libraries and JetBrains IDE support.
Flutter receives an Ecosystem score of \emph{0.74}.
It provides an official JetBrains plugin, built-in Material and Cupertino components, and a mature package ecosystem, but pub.dev contains substantially fewer packages than npm according to the selected normalisation method.

All frameworks receive a Quality Assurance score of \emph{1.00}.
They provide automated testing, static analysis tooling, and regular updates.
Flutter provides first-party testing and Dart analysis tooling, React Native uses Jest together with ESLint and TypeScript, and Ionic can use the testing and linting tooling of its selected frontend framework.

Overall, React Native achieves the highest result in the multi-criteria analysis.
Ionic and Flutter remain a close alternative with strong documentation, built-in UI functionality, and a mature mobile ecosystem.

\subsection{Prototyping}
Both Flutter and React Native were capable of implementing the required prototype functionality, including user interfaces, local storage, camera integration, and backend communication.

Flutter provided a structured development experience once the initial setup was completed.
Its widget-based approach and built-in components, such as \texttt{Scaffold}, \texttt{RefreshIndicator}, and \texttt{Snackbar}, made it possible to develop screens quickly.
Packages for shared preferences, camera access, local file paths, and sharing were also relatively easy to integrate.
However, configuring the Flutter, Dart, and Android environments required effort, and implementing a custom interface beyond the default Material Design appearance could require substantial widget-specific styling code.
The Android emulator was also resource-intensive and occasionally behaved differently from a physical device.

React Native provided a productive workflow for building and styling the interface.
Its component-based structure, CSS-like styling, and fast refresh functionality made it easy to create and adjust screens iteratively.
Navigation, SQLite storage, and camera functionality were implemented without major problems.
However, Android environment configuration required setting variables such as \texttt{ANDROID\_HOME} and \texttt{JAVA\_HOME}, and Java had to be downgraded to version 17 for Gradle to build the application.
Expo Go also did not support all required native functionality, meaning that the prototype had to be tested using a native Android build and USB debugging.
Content sharing required additional investigation because Expo-specific and general React Native packages supported different features.

Considering both the multi-criteria analysis and the prototyping results, React Native is the most suitable frontend framework for this project.
It achieves the highest overall score and aligns well with the team's existing JavaScript experience.
Its large ecosystem, strong community, familiar component model, straightforward styling, and fast refresh workflow outweigh the additional native Android and Expo-related configuration effort.

Flutter remains a strong alternative, particularly because of its structured widget model, mature documentation, and broad selection of built-in mobile functionality.
However, React Native is preferred because it better matches the team's existing skills and provided an efficient experience during prototyping.
